% Generated by publication/build_sections.py; do not edit by hand.
\paragraph{Informal verdict.}
\textbf{A --- Complete and apparently correct}, with \textbf{high confidence}.

\paragraph{Lean coverage (partial).}
Lean proves the decisive arbitrary-dimensional counterexample: for every finite-dimensional normed real space there is an integrable law whose normalized optimal quantization error has liminf zero and limsup infinity. The positive Zador asymptotics, universal lower bound, and unit-cube constant remain to be formalized.

\subsection{Audited informal answer}
Q838 --- Complete answer
\begin{enumerate}
\item Faithful restatement and source audit
\end{enumerate}
Let \((\Omega,\mathcal F,\mathbb P)\) be a probability space, let \(X:\Omega\to\mathbb R^n\) be a random vector, and fix a norm \(\|\cdot\|\) on \(\mathbb R^n\). For a finite set \(\Sigma\subset\mathbb R^n\), write

\[
d_\Sigma(x)=\inf_{a\in\Sigma}\|x-a\|,
\]

and define the \(k\)-point mean quantization error

\[
m_k(X)
   =\inf_{\#\Sigma=k}\mathbb E\,d_\Sigma(X).
\]

The question asks for general asymptotic results on \(m_k\), and in particular whether, for every integrable \(X\), the sequence \(k^{1/n}m_k(X)\) must have a limit.
The PDF also says that for the uniform distribution on \([-1,1]^n\) with the Euclidean norm, an earlier exercise proved

\[
k^{1/n}m_k\longrightarrow 2^{-n}.
\]

That assertion is indeed printed in the source. Luc Pommeret
As printed, it is false for every \(n\ge 2\). The correct limit has the form \(2q_{n,\|\cdot\|}\), where \(q_{n,\|\cdot\|}\) is the norm-dependent quantization constant defined below. For \(n=1\), \(q_{1,|\cdot|}=1/4\), so \(2q_{1,|\cdot|}=1/2=2^{-1}\); this probably explains why the erroneous expression is not detected in dimension one. The exact intended correction cannot be inferred from the PDF alone.
There are no dependence or independence assumptions: only one random vector is involved, and \(m_k\) depends solely on its law.

\begin{enumerate}
\item Main theorem
\end{enumerate}
Let \(\mu=\mathcal L(X)\), and write its Lebesgue decomposition as

\[
\mu=f(x)\,dx+\mu_s,
\]

where \(f\,dx\) is the absolutely continuous part and \(\mu_s\) is singular.
Put

\[
p=\frac{n}{n+1}.
\]

Let \(U\) be the uniform probability measure on the unit cube \(C=[0,1]^n\), and define

\[
q_{n,\|\cdot\|}
   :=\lim_{k\to\infty}
        k^{1/n}\inf_{\#\Sigma\le k}
        \int_C d_\Sigma(x)\,dx.
\]

The existence and positivity of this limit will be proved below.
Theorem
For every norm on \(\mathbb R^n\), the following statements hold.


Universal lower bound. If \(X\) is integrable, then

\[
\boxed{
\liminf_{k\to\infty} k^{1/n}m_k(X)
\ge
q_{n,\|\cdot\|}
\left(
   \int_{\mathbb R^n} f(x)^{\,n/(n+1)}\,dx
\right)^{(n+1)/n}.
}
\]

The right-hand side is interpreted as \(+\infty\) when the integral diverges.


Zador asymptotic under a higher moment. If, for some \(\eta>0\),

\[
\mathbb E\|X\|^{1+\eta}<\infty,
\]

then the limit exists and

\[
\boxed{
\lim_{k\to\infty} k^{1/n}m_k(X)
=
q_{n,\|\cdot\|}
\left(
   \int_{\mathbb R^n} f(x)^{\,n/(n+1)}\,dx
\right)^{(n+1)/n}.
}
\tag{2.1}
\]

In particular, the singular part of the law does not contribute to the leading \(k^{-1/n}\) term.


Integrability alone is not sufficient. For every \(n\ge1\) and every norm on \(\mathbb R^n\), there is an integrable discrete random vector \(X\) such that

\[
\boxed{
\liminf_{k\to\infty}k^{1/n}m_k(X)=0,
\qquad
\limsup_{k\to\infty}k^{1/n}m_k(X)=+\infty.
}
\tag{2.2}
\]



Thus the final question in Q838 has a definitive negative answer, even if limits in the extended real line are allowed.
Formula (2.1) is the \(r=1\) form of the classical Zador quantization theorem. The Euclidean \(r\)-power version under an \(r+\varepsilon\) moment condition, including the contribution of the absolutely continuous and singular parts, appears explicitly in Bucklew and Wise's theorem. Minds@UW+1
A self-contained proof of the form needed here follows.

\begin{enumerate}
\item Preliminary facts
\end{enumerate}
For a finite Borel measure \(\nu\) with finite first moment, define

\[
E_k(\nu)
   :=\inf_{\#\Sigma\le k}
      \int_{\mathbb R^n}d_\Sigma(x)\,d\nu(x).
\]

For a probability measure \(\mu=\mathcal L(X)\), this is \(m_k(X)\).
The infimum over \(\#\Sigma=k\) equals the infimum over \(\#\Sigma\le k\): a set with fewer than \(k\) points can be enlarged by adding arbitrary new points, and doing so cannot increase the distance.
We shall repeatedly use the following elementary properties.
Measure and spatial scaling
For \(c,t\ge0\),

\[
E_k(c\nu)=cE_k(\nu).
\]

If \(T(x)=a+tx\) and \(T_\#\nu\) is the pushforward, then

\[
E_k(T_\#\nu)=tE_k(\nu).
\]

Union of codebooks
For finite measures \(\nu,\rho\) and positive integers \(k,\ell\),

\[
E_{k+\ell}(\nu+\rho)
\le E_k(\nu)+E_\ell(\rho).
\tag{3.1}
\]

Indeed, take the union of a \(k\)-point codebook for \(\nu\) and an \(\ell\)-point codebook for \(\rho\).
Monotonicity under restriction
If \(0\le\nu\le\rho\), then

\[
E_k(\nu)\le E_k(\rho).
\tag{3.2}
\]

Stability under transport
The function \(x\mapsto d_\Sigma(x)\) is \(1\)-Lipschitz. Hence, if \(\nu,\rho\) have the same total mass and \(\pi\) is any coupling of them,

\[
\left\lvert 
\int d_\Sigma\,d\nu-\int d_\Sigma\,d\rho
\right\rvert 
\le
\int\|x-y\|\,d\pi(x,y).
\]

Taking infima over \(\Sigma\) gives

\[
|E_k(\nu)-E_k(\rho)|
\le
\inf_\pi\int\|x-y\|\,d\pi(x,y).
\tag{3.3}
\]

No optimal codebook is needed in the argument; arbitrarily near-optimal ones suffice.

\begin{enumerate}
\item The uniform-cube constant
\end{enumerate}
Let

\[
a_k:=E_k(\mathbf 1_C\,dx),
\qquad C=[0,1]^n.
\]

Since \(C\) has volume \(1\), \(\mathbf 1_C\,dx\) is the uniform probability measure.
4.1 Existence of the limit
Partition \(C\) into \(m^n\) congruent subcubes of side \(1/m\). In each subcube place a scaled copy of an \(N\)-point codebook for \(C\). Scaling gives

\[
a_{m^nN}\le \frac1m a_N.
\tag{4.1}
\]

Define

\[
q:=\inf_{N\ge1}N^{1/n}a_N.
\]

Clearly,

\[
k^{1/n}a_k\ge q.
\]

Conversely, fix \(N\) and put

\[
m=\left\lfloor (k/N)^{1/n}\right\rfloor.
\]

Then \(m^nN\le k\), so by monotonicity and (4.1),

\[
k^{1/n}a_k
\le
\frac{k^{1/n}}m\,a_N.
\]

As \(k\to\infty\),

\[
\frac{k^{1/n}}m\longrightarrow N^{1/n}.
\]

Consequently,

\[
\limsup_{k\to\infty}k^{1/n}a_k
\le N^{1/n}a_N.
\]

Taking the infimum over \(N\) gives

\[
\limsup_{k\to\infty}k^{1/n}a_k\le q.
\]

Therefore

\[
\boxed{
\lim_{k\to\infty}k^{1/n}a_k=q_{n,\|\cdot\|}.
}
\tag{4.2}
\]

4.2 Positivity and elementary bounds
Let

\[
\beta=\lambda_n\{x:\|x\|\le1\}
\]

be the Lebesgue volume of the norm unit ball.
If \(A\subset\mathbb R^n\) is measurable with volume \(v\), then for every \(a\in\mathbb R^n\),

\[
\begin{aligned}
\int_A\|x-a\|\,dx
&=
\int_0^\infty
   \lambda_n\bigl(A\setminus B(a,t)\bigr)\,dt\\
&\ge
\int_0^{(v/\beta)^{1/n}}
   \bigl(v-\beta t^n\bigr)\,dt\\
&=
\frac{n}{n+1}\beta^{-1/n}v^{1+1/n}.
\end{aligned}
\tag{4.3}
\]

For a \(k\)-point codebook, partition \(C\) into its Voronoi cells, whose volumes are \(v_1,\dots,v_r\), with \(r\le k\). Applying (4.3) and convexity gives

\[
\begin{aligned}
\int_Cd_\Sigma(x)\,dx
&\ge
\frac{n}{n+1}\beta^{-1/n}
\sum_{i=1}^r v_i^{1+1/n}\\
&\ge
\frac{n}{n+1}\beta^{-1/n}k^{-1/n}.
\end{aligned}
\]

Thus

\[
\boxed{
q_{n,\|\cdot\|}
\ge
\frac{n}{n+1}\beta^{-1/n}>0.
}
\tag{4.4}
\]

A regular cubical grid gives

\[
\boxed{
q_{n,\|\cdot\|}
\le
\int_{[-1/2,1/2]^n}\|x\|\,dx<\infty.
}
\tag{4.5}
\]


\begin{enumerate}
\item Compactly supported laws
\end{enumerate}
We first prove the asymptotic formula for compactly supported measures.
5.1 Piecewise constant densities on many small cubes
Put

\[
p=\frac{n}{n+1};
\qquad
\frac1p=1+\frac1n.
\]

Consider a density

\[
g_k=\sum_{i=1}^{M_k}c_{i,k}\mathbf 1_{Q_{i,k}},
\]

where the \(Q_{i,k}\) are cubes of a common side length \(h_k\), contained in a fixed bounded region. Assume

\[
L_k:=k^{1/n}h_k\longrightarrow\infty.
\tag{5.1}
\]

Then

\[
M_k=O(h_k^{-n})=O(k/L_k^n)=o(k).
\tag{5.2}
\]

Set

\[
s_{i,k}=c_{i,k}^{\,p}|Q_{i,k}|,
\qquad
S_k=\sum_i s_{i,k}=\int g_k^p,
\]

and

\[
A_{i,k}=c_{i,k}|Q_{i,k}|^{1+1/n}.
\]

Since \(p(1+1/n)=1\),

\[
A_{i,k}^{\,p}=s_{i,k}.
\tag{5.3}
\]

We claim that, provided the total masses remain bounded,

\[
\boxed{
k^{1/n}E_k(g_k\,dx)
-
q_{n,\|\cdot\|}S_k^{1/p}
\longrightarrow0.
}
\tag{5.4}
\]

Upper bound
Fix \(\varepsilon>0\). By (4.2), there is an integer \(R\) such that

\[
a_r\le(q+\varepsilon)r^{-1/n}
\qquad(r\ge R).
\tag{5.5}
\]

If \(S_k>0\), put

\[
w_{i,k}=\frac{s_{i,k}}{S_k}.
\]

Give \(R\) points to every positive-density cube, and distribute the remaining points proportionally to \(w_{i,k}\). More precisely, with

\[
K_k=k-RM_k,
\]

set

\[
r_{i,k}=R+\lfloor K_kw_{i,k}\rfloor.
\]

By (5.2), \(K_k/k\to1\), and \(\sum_i r_{i,k}\le k\). Moreover,

\[
r_{i,k}\ge K_kw_{i,k}.
\]

Using scaled copies of unit-cube codebooks,

\[
E_k(g_kdx)
\le
\sum_i A_{i,k}a_{r_{i,k}}.
\]

Consequently,

\[
\begin{aligned}
E_k(g_kdx)
&\le
(q+\varepsilon)
\sum_iA_{i,k}r_{i,k}^{-1/n}\\
&\le
(q+\varepsilon)K_k^{-1/n}
\sum_iA_{i,k}w_{i,k}^{-1/n}.
\end{aligned}
\]

From (5.3),

\[
A_{i,k}
=(S_kw_{i,k})^{1/p},
\]

and hence

\[
A_{i,k}w_{i,k}^{-1/n}
=
S_k^{1/p}w_{i,k}.
\]

Therefore

\[
\limsup_{k\to\infty}
k^{1/n}E_k(g_kdx)
\le
(q+\varepsilon)S_k^{1/p}+o(1).
\tag{5.6}
\]

Lower bound: the boundary firewall
Let \(\Sigma\) be a global codebook with at most \(k\) points. Assign every point of \(\Sigma\) to at most one cube.
For every cube \(Q_{i,k}\), place a finite \(\delta_k\)-net on its boundary. For \(n\ge2\), such a net can be chosen with at most

\[
C\left(\frac{h_k}{\delta_k}\right)^{n-1}
\]

points. For \(n=1\), the two endpoints suffice.
Choose, for \(n\ge2\),

\[
\ell_k=L_k^{1/(2(n-1))},
\qquad
\delta_k=\frac{k^{-1/n}}{\ell_k}.
\]

Then

\[
k^{1/n}\delta_k=\ell_k^{-1}\longrightarrow0,
\tag{5.7}
\]

while the total number of firewall points is

\[
O\left(
M_k\left(\frac{h_k}{\delta_k}\right)^{n-1}
\right)
=
O\left(
\frac{k}{L_k^n}(L_k\ell_k)^{n-1}
\right)
=
O(kL_k^{-1/2})
=o(k).
\tag{5.8}
\]

If a nearest point of the global codebook to \(x\in Q_{i,k}\) is not assigned to \(Q_{i,k}\), the line segment joining \(x\) to that point meets \(\partial Q_{i,k}\). A firewall point lies within \(C\delta_k\) of the intersection point. Therefore, if \(\Gamma_{i,k}\) consists of the codepoints assigned to \(Q_{i,k}\) together with its firewall,

\[
d_\Sigma(x)
\ge d_{\Gamma_{i,k}}(x)-C\delta_k
\qquad(x\in Q_{i,k}).
\tag{5.9}
\]

Let \(r_{i,k}\) be an upper bound for \(\#\Gamma_{i,k}\). By (5.8),

\[
\sum_i r_{i,k}\le k+o(k).
\tag{5.10}
\]

Scaling the unit-cube problem and using the defining inequality

\[
a_r\ge qr^{-1/n},
\]

we get

\[
E_k(g_kdx)
\ge
q\sum_iA_{i,k}r_{i,k}^{-1/n}
-o(k^{-1/n}).
\tag{5.11}
\]

For positive \(A_i,r_i\), Hölder's inequality gives

\[
\sum_i A_i r_i^{-1/n}
\ge
\left(\sum_iA_i^p\right)^{1/p}
\left(\sum_i r_i\right)^{-1/n}.
\tag{5.12}
\]

Indeed,

\[
\sum_iA_i^p
=
\sum_i
   (A_ir_i^{-1/n})^p r_i^{\,1-p}
\le
\left(\sum_iA_ir_i^{-1/n}\right)^p
\left(\sum_i r_i\right)^{1-p}.
\]

Using (5.3), (5.10), and (5.12) in (5.11),

\[
\liminf_{k\to\infty}
k^{1/n}E_k(g_kdx)
\ge qS_k^{1/p}+o(1).
\tag{5.13}
\]

Together, (5.6) and (5.13) prove (5.4).

5.2 Arbitrary compactly supported densities
Let \(f\ge0\) be integrable and supported in a fixed bounded cube.
For a grid of mesh \(h\), let \(f_h\) be the function which, on each grid cube \(Q\), is equal to the average of \(f\) on \(Q\):

\[
f_h|_Q=\frac1{|Q|}\int_Qf.
\]

Then

\[
\|f_h-f\|_{L^1}\longrightarrow0
\qquad(h\downarrow0).
\tag{5.14}
\]

This follows first for continuous functions by uniform continuity and then for \(L^1\) functions by approximation, since the averaging operators are \(L^1\)-contractions.
Write

\[
\varepsilon(h)=\|f_h-f\|_1.
\]

Choose \(h_k\downarrow0\) so slowly that

\[
L_k:=k^{1/n}h_k\longrightarrow\infty
\quad\text{and}\quad
L_k\varepsilon(h_k)\longrightarrow0.
\tag{5.15}
\]

Such a choice is always possible. For example, choose numbers \(\delta_j\downarrow0\) such that \(\varepsilon(h)\le j^{-2}\) for \(h\le\delta_j\), and let \(L_k\) be the largest \(j\) such that

\[
jk^{-1/n}\le \min(\delta_j,j^{-1}).
\]

Then set \(h_k=L_kk^{-1/n}\).
On each grid cube, the measures \(f\,dx\) and \(f_{h_k}\,dx\) have the same total mass. Couple their common part identically, and couple the remaining equal masses within the same cube. Since the diameter of a grid cube is at most \(Ch_k\),

\[
W_1(f\,dx,f_{h_k}\,dx)
\le
Ch_k\|f-f_{h_k}\|_1.
\]

By (3.3) and (5.15),

\[
k^{1/n}
\left\lvert 
E_k(f\,dx)-E_k(f_{h_k}\,dx)
\right\rvert 
\longrightarrow0.
\tag{5.16}
\]

Furthermore, for \(0<p<1\),

\[
|u^p-v^p|\le |u-v|^p,
\]

and on a set of finite volume,

\[
\int|f-f_{h_k}|^p
\le
C\|f-f_{h_k}\|_1^p.
\]

Thus

\[
\int f_{h_k}^p\longrightarrow\int f^p.
\tag{5.17}
\]

Applying (5.4) to \(f_{h_k}\) and using (5.16)--(5.17) yields

\[
\boxed{
\lim_{k\to\infty}
k^{1/n}E_k(f\,dx)
=
q
\left(\int f^p\right)^{1/p}.
}
\tag{5.18}
\]


5.3 Compactly supported singular measures
First suppose that a finite measure \(\nu\) is supported by a compact Lebesgue-null set \(K\).
Let \(N_K(h)\) denote the number of grid cubes of side \(h\) meeting \(K\). Their union lies in a \(Ch\)-neighborhood \(K^{Ch}\) of \(K\), so

\[
N_K(h)h^n\le \lambda_n(K^{Ch}).
\]

Since \(K\) is compact and null,

\[
\lambda_n(K^{Ch})\longrightarrow0.
\]

Given \(\varepsilon>0\), put

\[
h_k=\left(\frac{2\varepsilon}{k}\right)^{1/n}.
\]

For all sufficiently large \(k\),

\[
N_K(h_k)\le \frac{k}{2}.
\]

Placing one point in every cube meeting \(K\) gives covering radius at most \(Ch_k\). Therefore

\[
E_k(\nu)\le C\nu(\mathbb R^n)h_k,
\]

and consequently

\[
\limsup_{k\to\infty}k^{1/n}E_k(\nu)
\le
C\nu(\mathbb R^n)\varepsilon^{1/n}.
\]

Letting \(\varepsilon\downarrow0\),

\[
\boxed{
k^{1/n}E_k(\nu)\longrightarrow0.
}
\tag{5.19}
\]

Now let \(\nu\) be any singular measure supported in a fixed compact set \(K_0\). There is a null Borel set \(N\) carrying all its mass. By inner regularity, choose increasing compact sets \(K_j\subset N\) such that

\[
\nu(K_0\setminus K_j)\longrightarrow0.
\]

The restriction \(\nu|_{K_j}\) satisfies (5.19).
For every finite measure \(\rho\) supported in \(K_0\), a regular grid gives a universal estimate

\[
E_\ell(\rho)
\le
C_{K_0}\rho(\mathbb R^n)\ell^{-1/n}.
\tag{5.20}
\]

Allocate \((1-\varepsilon)k\) points to \(\nu|_{K_j}\) and \(\varepsilon k\) points to the remaining measure. From (3.1), (5.19), and (5.20),

\[
\limsup_{k\to\infty}k^{1/n}E_k(\nu)
\le
C_{K_0}
\nu(K_0\setminus K_j)\varepsilon^{-1/n}.
\]

For instance, choose

\[
\varepsilon
=
\nu(K_0\setminus K_j)^{\,n/(n+1)}.
\]

The right-hand side then tends to zero as \(j\to\infty\). Hence every compactly supported singular measure satisfies (5.19).

5.4 The compact-support formula
Let

\[
\mu=f\,dx+\mu_s
\]

be compactly supported.
By monotonicity,

\[
E_k(\mu)\ge E_k(f\,dx),
\]

so (5.18) gives the lower bound.
For the upper bound, allocate \((1-\varepsilon)k\) points to \(f\,dx\) and \(\varepsilon k\) points to \(\mu_s\). Equations (3.1), (5.18), and (5.19) give

\[
\limsup_{k\to\infty}k^{1/n}E_k(\mu)
\le
(1-\varepsilon)^{-1/n}
q\left(\int f^p\right)^{1/p}.
\]

Letting \(\varepsilon\downarrow0\),

\[
\boxed{
\lim_{k\to\infty}k^{1/n}E_k(\mu)
=
q\left(\int f^p\right)^{1/p}
}
\tag{5.21}
\]

for every compactly supported finite measure.

\begin{enumerate}
\item Unbounded laws and the moment condition
\end{enumerate}
We need a uniform tail estimate.
6.1 A Pierce-type bound
Let \(s>1\). There is a constant \(C=C(n,\|\cdot\|,s)\) such that, for every finite measure \(\nu\),

\[
\boxed{
E_k(\nu)
\le
Ck^{-1/n}
\nu(\mathbb R^n)^{\,1-1/s}
\left(\int\|x\|^s\,d\nu(x)\right)^{1/s}.
}
\tag{6.1}
\]

Proof
It suffices first to consider a probability measure \(P\) satisfying

\[
\int\|x\|^s\,dP(x)\le1.
\]

Partition space into

\[
A_0=\{\|x\|\le1\},
\qquad
A_j=\{2^{j-1}<\|x\|\le2^j\},
\quad j\ge1.
\]

Put

\[
p_j=P(A_j),
\qquad
R_0=1,\quad R_j=2^j,
\qquad
a_j=p_jR_j.
\]

Let again \(p=n/(n+1)\). We claim

\[
S:=\sum_{j\ge0}a_j^p\le C.
\tag{6.2}
\]

For \(j\ge1\), put \(b_j=p_jR_j^s\). Then \(\sum_jb_j\le C_s\), and

\[
a_j^p=b_j^pR_j^{-p(s-1)}.
\]

Since \(1-p=1/(n+1)\), write

\[
R_j^{-p(s-1)}
=
\bigl(R_j^{-n(s-1)}\bigr)^{1-p}.
\]

Hölder's inequality gives

\[
\sum_{j\ge1}a_j^p
\le
\left(\sum_jb_j\right)^p
\left(\sum_jR_j^{-n(s-1)}\right)^{1-p}<\infty,
\]

proving (6.2).
Set

\[
w_j=\frac{a_j^p}{S}.
\]

Use one codepoint at the origin. For every \(j\) with \((k-1)w_j\ge1\), allocate

\[
\ell_j=\lfloor(k-1)w_j\rfloor
\]

additional codepoints to the ball of radius \(R_j\). Such a ball can be covered with \(\ell_j\) points and covering radius at most

\[
CR_j\ell_j^{-1/n}.
\]

For the remaining shells, use the origin.
On the allocated shells,

\[
\begin{aligned}
\sum_{\text{allocated}}p_jCR_j\ell_j^{-1/n}
&\le
Ck^{-1/n}
\sum_j a_jw_j^{-1/n}\\
&=
Ck^{-1/n}S^{1/p}
\sum_jw_j\\
&\le Ck^{-1/n}.
\end{aligned}
\]

On the unallocated shells, \(w_j<(k-1)^{-1}\), and therefore

\[
w_j^{1/p}
=w_j^{1+1/n}
\le
(k-1)^{-1/n}w_j.
\]

Hence

\[
\sum_{\text{unallocated}}a_j
=
S^{1/p}\sum_{\text{unallocated}}w_j^{1/p}
\le Ck^{-1/n}.
\]

This proves \(E_k(P)\le Ck^{-1/n}\) in the normalized case.
Scaling space by

\[
\sigma=\left(\int\|x\|^s\,dP\right)^{1/s}
\]

gives the probability version of (6.1). Normalizing a finite measure by its total mass gives the displayed finite-measure estimate.

6.2 Proof of the full asymptotic formula
Assume now

\[
\int\|x\|^{1+\eta}\,d\mu(x)<\infty.
\]

Set \(s=1+\eta\).
First, the density term is finite. Indeed, with

\[
w(x)=1+\|x\|^s,
\]

Hölder's inequality gives

\[
\begin{aligned}
\int f^p
&=
\int(fw)^p w^{-p}\\
&\le
\left(\int fw\right)^p
\left(
   \int w^{-p/(1-p)}
\right)^{1-p}.
\end{aligned}
\]

Here

\[
\frac{p}{1-p}=n,
\]

and

\[
\int_{\mathbb R^n}(1+\|x\|^s)^{-n}\,dx<\infty
\]

because \(s>1\).
Let

\[
\mu_R=\mu|_{\{\|x\|\le R\}},
\qquad
\tau_R=\mu|_{\{\|x\|>R\}}.
\]

By monotonicity and the compact-support formula,

\[
\begin{aligned}
\liminf_{k\to\infty}k^{1/n}E_k(\mu)
&\ge
q
\left(
  \int_{\{\|x\|\le R\}} f^p
\right)^{1/p}.
\end{aligned}
\]

Letting \(R\to\infty\) and using monotone convergence,

\[
\liminf_{k\to\infty}k^{1/n}E_k(\mu)
\ge
q\left(\int f^p\right)^{1/p}.
\tag{6.3}
\]

For the upper bound, fix \(0<\varepsilon<1\). Allocate

\[
k_1=\lfloor(1-\varepsilon)k\rfloor
\]

points to \(\mu_R\), and the remaining \(k_2\) points to \(\tau_R\). From (3.1),

\[
E_k(\mu)\le E_{k_1}(\mu_R)+E_{k_2}(\tau_R).
\]

The compact-support formula gives

\[
\lim_{k\to\infty}
k^{1/n}E_{k_1}(\mu_R)
=
(1-\varepsilon)^{-1/n}
q\left(
  \int_{\{\|x\|\le R\}}f^p
\right)^{1/p}.
\]

By (6.1),

\[
k^{1/n}E_{k_2}(\tau_R)
\le
C\varepsilon^{-1/n}
\left(
   \int_{\{\|x\|>R\}}\|x\|^s\,d\mu(x)
\right)^{1/s}.
\]

The last expression tends to zero as \(R\to\infty\). Thus

\[
\limsup_{k\to\infty}k^{1/n}E_k(\mu)
\le
(1-\varepsilon)^{-1/n}
q\left(\int f^p\right)^{1/p}.
\]

Finally let \(\varepsilon\downarrow0\). Together with (6.3), this proves (2.1).
The same argument also proves the universal lower bound without any moment beyond what is needed to make \(m_k\) finite: apply the compact-support result to \(\mu_R\) and let \(R\to\infty\).

\begin{enumerate}
\item The uniform cube and the error in the PDF
\end{enumerate}
Let \(X\) be uniform on \([-1,1]^n\). Its density is

\[
f(x)=2^{-n}\mathbf 1_{[-1,1]^n}(x).
\]

Therefore

\[
\begin{aligned}
\int f(x)^p\,dx
&=
2^n(2^{-n})^p\\
&=
2^{n(1-p)}
=
2^{n/(n+1)}.
\end{aligned}
\]

Raising to the power \(1/p=(n+1)/n\) gives

\[
\left(\int f^p\right)^{1/p}=2.
\]

Hence the correct answer is

\[
\boxed{
\lim_{k\to\infty}k^{1/n}m_k
=
2q_{n,\|\cdot\|}.
}
\tag{7.1}
\]

For the Euclidean norm, \(\beta=\omega_n\), the volume of the Euclidean unit ball. From (4.4),

\[
2q_{n,2}
\ge
\frac{2n}{n+1}\omega_n^{-1/n}.
\tag{7.2}
\]

For example, when \(n=2\),

\[
2q_{2,2}
\ge
\frac{4}{3\sqrt{\pi}}
%
\frac14=2^{-2}.
\]

Thus the printed value \(2^{-n}\) is impossible already in dimension two.

\begin{enumerate}
\item Counterexample under mere integrability
\end{enumerate}
We now prove the strong nonconvergence assertion (2.2).
Fix \(n\ge1\), fix any norm on \(\mathbb R^n\), and choose \(u\in\mathbb R^n\) such that

\[
\|u\|=1.
\]

8.1 A decreasing summable block sequence
Define recursively

\[
N_1=1,
\]

and, for \(s\ge1\),

\[
L_s=2^{2ns}N_s^2,
\qquad
N_{s+1}=N_s+L_s.
\]

Set

\[
M_s=2^{-s}N_s^{-1/n},
\qquad
h_s=\frac{M_s}{L_s}.
\]

For every integer \(j\) in the \(s\)-th block,

\[
N_s\le j<N_{s+1},
\]

define

\[
a_j=h_s.
\]

Then

\[
\sum_{j\ge1}a_j
=
\sum_{s\ge1}L_sh_s
=
\sum_{s\ge1}M_s
\le
\sum_{s\ge1}2^{-s}
=1.
\tag{8.1}
\]

Moreover,

\[
h_s
=
2^{-(2n+1)s}N_s^{-(2+1/n)},
\]

so \((a_j)\) is nonincreasing.
Put

\[
A_k=\sum_{j\ge k}a_j.
\]

8.2 Definition of the random vector
Let

\[
y_j=4^j u,
\qquad
p_j=\frac{a_j}{4^j}.
\]

Since \(\sum_jp_j\le \frac14\sum_ja_j<1\), define \(X\) by

\[
\mathbb P(X=y_j)=p_j,
\qquad
\mathbb P(X=0)=1-\sum_{j\ge1}p_j.
\]

By Tonelli's theorem,

\[
\mathbb E\|X\|
=
\sum_{j\ge1}p_j\|y_j\|
=
\sum_{j\ge1}a_j
\le1.
\tag{8.2}
\]

Thus \(X\) is integrable.
8.3 Two-sided estimate for \(m_k\)
For the upper bound, use

\[
\Sigma_k=\{0,y_1,\dots,y_{k-1}\}.
\]

Then

\[
m_k
\le
\sum_{j\ge k}p_j\|y_j\|
=
A_k.
\tag{8.3}
\]

For the lower bound, consider the closed norm-balls

\[
B_j=\overline B\left(y_j,\frac{4^j}{4}\right).
\]

They are pairwise disjoint. Indeed, for \(i<j\),

\[
\|y_j-y_i\|=4^j-4^i
%
\frac{4^j+4^i}{4}.
\]

A \(k\)-point codebook can meet at most \(k\) of these balls. If \(B_j\) contains no codepoint, then

\[
d_\Sigma(y_j)\ge\frac{4^j}{4}.
\]

Hence

\[
\int d_\Sigma\,d\mu
\ge
\frac14\sum_{\substack{j:\,B_j\cap\Sigma=\varnothing}}a_j.
\]

Since \((a_j)\) is nonincreasing, deleting at most \(k\) terms leaves sum at least \(\sum_{j\ge k+1}a_j\). Therefore

\[
m_k\ge\frac14A_{k+1}.
\tag{8.4}
\]

Combining (8.3) and (8.4),

\[
\boxed{
\frac14A_{k+1}\le m_k\le A_k.
}
\tag{8.5}
\]

8.4 A subsequence tending to zero
At the beginning of the \(s\)-th block,

\[
\begin{aligned}
A_{N_s}
&=
\sum_{t\ge s}M_t\\
&=
\sum_{t\ge s}2^{-t}N_t^{-1/n}\\
&\le
N_s^{-1/n}\sum_{t\ge s}2^{-t}\\
&=
2^{1-s}N_s^{-1/n}.
\end{aligned}
\]

Thus, by (8.3),

\[
N_s^{1/n}m_{N_s}
\le2^{1-s}\longrightarrow0.
\tag{8.6}
\]

8.5 A subsequence tending to infinity
Put

\[
K_s=N_s+\left\lfloor\frac{L_s}{2}\right\rfloor.
\]

At least \(L_s/4\) terms of the \(s\)-th block occur after \(K_s\), so

\[
A_{K_s+1}
\ge
\frac{L_s}{4}h_s
=
\frac{M_s}{4}.
\]

By (8.4),

\[
m_{K_s}\ge\frac{M_s}{16}.
\]

Also \(K_s\ge L_s/2\), hence

\[
\begin{aligned}
K_s^{1/n}m_{K_s}
&\ge
\frac1{16}
\left(\frac{L_s}{2}\right)^{1/n}
2^{-s}N_s^{-1/n}\\
&=
\frac{2^{-1/n}}{16}
\cdot
2^{-s}
\cdot
2^{2s}
N_s^{2/n}
N_s^{-1/n}\\
&=
\frac{2^{-1/n}}{16}
\,2^sN_s^{1/n}
\longrightarrow+\infty.
\end{aligned}
\tag{8.7}
\]

Equations (8.6) and (8.7) prove

\[
\liminf_{k\to\infty}k^{1/n}m_k=0,
\qquad
\limsup_{k\to\infty}k^{1/n}m_k=+\infty.
\]

This completes the counterexample.

\begin{enumerate}
\item Small and boundary cases
\end{enumerate}
9.1 Dimension one
For the usual absolute value on \(\mathbb R\),

\[
q_{1,|\cdot|}=\frac14.
\]

Indeed, if an interval of length \(\ell\) is represented by one point, the minimum possible integral distance is \(\ell^2/4\). If the Voronoi intervals have lengths \(\ell_1,\dots,\ell_k\) with total length \(1\),

\[
\int_0^1d_\Sigma(x)\,dx
\ge
\frac14\sum_{i=1}^k\ell_i^2
\ge
\frac1{4k}.
\]

Equality is attained by the midpoints of \(k\) equal subintervals.
Thus, under a \(1+\eta\) moment,

\[
\boxed{
\lim_{k\to\infty}k\,m_k(X)
=
\frac14
\left(\int_{\mathbb R}\sqrt{f(x)}\,dx\right)^2.
}
\tag{9.1}
\]

For \(X\) uniform on \([-1,1]\), one has exactly

\[
m_k=\frac1{2k}.
\]

9.2 Finite support and atoms
If \(X\) has support of cardinality \(r\), then

\[
m_k=0\qquad(k\ge r).
\]

Its absolutely continuous density is zero, so formula (2.1) gives zero, as required.
More generally, atoms and compactly supported singular components have zero contribution at the \(k^{-1/n}\) scale.
9.3 Mixed distributions
If

\[
\mu=a\,\nu_{\mathrm{ac}}+(1-a)\nu_s
\]

and \(\nu_{\mathrm{ac}}\) has density \(g\), then \(f=ag\), and

\[
\left(\int (ag)^p\right)^{1/p}
=
a\left(\int g^p\right)^{1/p}.
\]

Thus the leading coefficient is linear in the mass \(a\) of the absolutely continuous component.
9.4 Necessity of integrability for finite error
If \(\mathbb E\|X\|=\infty\), then for every finite \(\Sigma\), with

\[
R=\max_{a\in\Sigma}\|a\|,
\]

the reverse triangle inequality gives

\[
d_\Sigma(X)\ge \|X\|-R.
\]

Hence \(\mathbb E d_\Sigma(X)=\infty\), and \(m_k=\infty\).
The \(1+\eta\) moment in the theorem is a sufficient tail condition, not a necessary one. The counterexample proves, however, that it cannot simply be replaced by the first-moment assumption for arbitrary laws.

\begin{enumerate}
\item Final conclusion
\end{enumerate}
The mathematically correct general result is

\[
\lim_{k\to\infty}k^{1/n}m_k
=
q_{n,\|\cdot\|}
\left(
\int_{\mathbb R^n}
f(x)^{n/(n+1)}\,dx
\right)^{(n+1)/n}
\]

whenever \(X\) has a moment of some order strictly greater than \(1\). The density \(f\) is that of the absolutely continuous part of the law; singular mass is asymptotically negligible at this scale.
For the uniform law on \([-1,1]^n\), the correct constant is

\[
2q_{n,\|\cdot\|},
\]

not \(2^{-n}\) except accidentally when \(n=1\).
Finally, mere integrability does not guarantee convergence: for every dimension and every norm there is an integrable discrete law for which

\[
\liminf k^{1/n}m_k=0
\quad\text{and}\quad
\limsup k^{1/n}m_k=+\infty.
\]
